% TOP_PROBLEM: {"problemId": "problem.skorobogatov-brauer-manin-conjecture-for-k3-surfaces", "canonicalTitle": "Brauer-Manin obstruction to the Hasse principle for K3 surfaces", "releaseRank": 383, "primaryDomain": "number_theory_arithmetic_geometry", "primaryDomainLabel": "Number theory & arithmetic geometry", "displayStatus": "open_with_solved_subcases", "releaseStatus": "open_with_solved_subcases"}
% !TeX program = lualatex
% Definition and sources only; this file is not an LLM solution attempt.
\documentclass[11pt]{article}
\usepackage[margin=25mm]{geometry}
\usepackage{fontspec}
\setmainfont{Latin Modern Roman}
\usepackage{amsmath,amssymb,mathtools}
\usepackage{unicode-math}
\setmathfont{Latin Modern Math}
\usepackage{xurl,hyperref}
\hypersetup{colorlinks=true,urlcolor=blue}
\setcounter{secnumdepth}{0}
\setlength{\parindent}{0pt}
\setlength{\parskip}{5pt}
\setlength{\emergencystretch}{3em}
\begin{document}
\section*{\texorpdfstring{Brauer-Manin obstruction to the Hasse principle for K3 surfaces}{Brauer-Manin obstruction to the Hasse principle for K3 surfaces}}\label{383}
\subsection{Definitions and mathematical statement}
A K3 surface over a number field $k$ is a smooth projective geometrically integral surface $X$ with $K_X\cong\mathcal O_X$ and $H^1(X,\mathcal O_X)=0$. Let ${\mathbb{A}}_k$ be the adele ring and $\operatorname{Br}(X)=H^2_{\mathrm{et}}(X,\mathbf G_m)$. Define
\[
 X({\mathbb{A}}_k)^{\operatorname{Br}(X)}=
 \left\{(x_v):\sum_v\operatorname{inv}_v(\alpha(x_v))=0
 \text{ for every }\alpha\in\operatorname{Br}(X)\right\},
\]
using the local invariant maps $\operatorname{Br}(k_v)\to{\mathbb{Q}}/{\mathbb{Z}}$; the sum has only finitely many nonzero terms for a fixed class and adelic point. The conjecture is that nonemptiness of this set implies $X(k)\ne\varnothing$. One adelic point must annihilate the entire Brauer group simultaneously. Density of rational points, weak approximation, and obstruction by one single Brauer class are not required.
\par\smallskip\textit{Formulation and concept references:} \href{https://arxiv.org/html/2502.08723v1}{[S1]}.

\subsection{Short English statement}
For every K3 surface over a number field, does surviving all Brauer--Manin tests guarantee at least one rational point?
\subsection{Sources}
\begin{itemize}
\item[S1]S. Monnet, Explicit bounds on the transcendental Brauer group of K3 surfaces with principal complex multiplication, arXiv:2502.08723v1 (2025); pairing convention from Tawfik-Newton arXiv:2311.00650v2.
\url{https://arxiv.org/html/2502.08723v1}.
\textit{Formulation source.}
\end{itemize}
\end{document}
